\subsection{Constructive ambiguity from two feasible pose designs}
Wide upper bounds alone cannot distinguish intrinsic ambiguity from a loose
audit. We therefore construct lower bounds from two allowed pose designs,
using the classical Gaussian testing argument \citep{donoho1994}. This is a
diagnostic on the same declared class, not a new minimax principle.

Let $A_0,A_1$ be two allowed observation operators, and let
$\rho_j$ belong to $\mathcal C=\{\rho:\norm{\rho-\rho_0^{\rm pilot}}\le B\}$.
Both designs have identity observation covariance. If
\begin{equation}
 D=\norm{A_1\rho_1-A_0\rho_0}<2\Phi^{-1}(1-\alpha),
 \label{eq:two-pose-distance}
\end{equation}
then every uniformly $1-\alpha$ valid interval whose length is deterministic
must have half-length at least
$|\langle\ell,\rho_1-\rho_0\rangle|/2$. Indeed, if a shorter interval
contained either feature with probability at least $1-\alpha$, its membership
test for one feature would distinguish the two hypotheses with both error
probabilities at most $\alpha$. The minimum sum of errors for their Gaussian
distributions is $2\Phi(-D/2)>2\alpha$, a contradiction. This conclusion does
not bound the expected length of arbitrary variable-length intervals.

Write $p=A_1\rho_0^{\rm pilot}-A_0\rho_0^{\rm pilot}$ and choose a distance
budget $\tau<2\Phi^{-1}(1-\alpha)$. For any observation weights $w$,
\begin{equation}
 G(w)=B\sum_{j=0}^1\norm{\ell-A_j^*w}-w^\top p+\tau\norm w
 \label{eq:two-pose-dual}
\end{equation}
upper-bounds the largest feasible feature gap for this fixed pair of pose
designs. To see this, insert $A_j^*w$ into the two feature pairings, apply
the density-ball bounds, and use
$w^\top(A_1\rho_1-A_0\rho_0)\le\tau\norm w$.

The same weights yield explicit continuous witnesses. Put
$h_j=\ell-A_j^*w$, choose $r_j\ge\norm{h_j}$, and set
\begin{equation}
 \widetilde\rho_1=\rho_0^{\rm pilot}+Bh_1/r_1,\qquad
 \widetilde\rho_0=\rho_0^{\rm pilot}-Bh_0/r_0.
\end{equation}
If $r_j=0$, its perturbation is set to zero. Their mean difference is
$p+\sum_j(B/r_j)(A_j\ell-A_jA_j^*w)$ and their feature difference is
$\sum_j(B/r_j)(\norm\ell^2-w^\top A_j\ell)$. Let $\overline D$ bound the
mean-distance norm. Multiplying both densities by
$\lambda=\min(1,\tau/\overline D)$ (one if $\overline D=0$) makes them feasible for
Equation~\eqref{eq:two-pose-distance}. This scaling preserves the original
class when $\norm{\rho_0^{\rm pilot}}\le B$, because zero is in the ball
and the ball is convex. Thus the scaled feature gap supplies a constructive
lower bound even if optimization is incomplete.

We minimize Equation~\eqref{eq:two-pose-dual} by iteratively reweighted
quadratic solves and retain both the best feasible pair and the best dual
weights, which may occur at different iterations. Continuous Gram integration
errors pad each residual norm and the mean-distance norm. Translation is a
real orthogonal rotation of each Fourier cosine/sine pair, so the phase-wrapped
Gram and target preserve those error norms. The witnesses are continuous
target/adjoint fields plus the known pilot; they are not restricted to a voxel
representation. Floating-point magnitude guards are additional numerical
checks, not validated interval arithmetic.

The prespecified development grid uses the existing 128-particle, radius-5
geometries and broad targets on three stacks. It includes identical nominal
poses, antipodal coherent rotations, and antipodal random joint-ball poses at
one and two degrees with $0.5\,\AA$ translation radius. The same random unit
vectors are used at both angular budgets. Independent conic programs check
small Hilbert-space cases; deliberately crude quadrature is checked against
the exact continuous sinc Gram; nonlinear constant-cell projections check
the phase convention. A serialization failure and the first complete grid
remain archived. The final rerun adds the numerical guards and separate
witness records requested by the focused mathematical audit.

\begin{table*}[t]
\centering
\begin{tabular}{llrrrrr}
\toprule
Stack & Target & Nominal & Coherent $1^\circ$ & Random $1^\circ$ & Coherent $2^\circ$ & Random $2^\circ$ \\
\midrule
10028 & center & 0.087 & 0.087 & 0.088 & 0.087 & 0.088 \\
10028 & contrast & 0.106 & 0.108 & 0.107 & 0.113 & 0.107 \\
10049 & center & 0.123 & 0.123 & 0.125 & 0.123 & 0.126 \\
10049 & contrast & 0.147 & 0.148 & 0.148 & 0.151 & 0.150 \\
10076 & center & 0.106 & 0.106 & 0.106 & 0.106 & 0.106 \\
10076 & contrast & 0.126 & 0.127 & 0.126 & 0.132 & 0.127 \\
\bottomrule
\end{tabular}
\caption{Constructive lower bounds on the half-width of any uniformly valid deterministic-length interval, divided by the no-data half-width $B\norm\ell$. Each entry uses two feasible antipodal pose configurations, continuous densities in the original ball and a Gaussian testing distance below $2\Phi^{-1}(0.95)$. Random poses are fixed design draws. All 30 fixed-pair modulus gaps are below 0.5\%; this does not certify maximization over poses. A small entry cannot establish that a larger upper bound is loose.}
\label{tab:two-pose-modulus}
\end{table*}

The fixed-pair problems are solved to relative brackets below 0.5\%. Their
lower bounds remain substantially smaller than several pose-aware upper
intervals. This leaves a gap: the selected configurations need not maximize
ambiguity, so these results do not prove that the larger upper bounds can be
attained or substantially reduced. The testing statement also retains the
declared white-Gaussian noise, density and nuisance assumptions.
